\section{服务范围}
\subsection{深度清洁工艺技术现状及适应性研究}
\begin{circlist}
\item 深度清洁工艺技术的主要方式；
\item 优缺点分析及海上现状适用性分析。
\end{circlist}

\subsection{井筒高效清洁技术配套理论研究}
\begin{circlist}
\item 砂砾、铁屑、结垢物、蜡等井筒内的运移及分布规律研究，并基于室内试验验证规律的合理性；
\item 高强磁性体及高温磁性体的优选，最大磁性布置规律及分布式磁性吸附防卡及延长寿命研究；
\item 水力脉冲解堵清洁技术研究；
\item 清洁工具关键部件对不同附着物清洁效果评价研究。
\end{circlist}

\subsection{井筒高效清洁技术配套工具设计}
\begin{circlist}
\item “四合一”通刮刷磁集成式井筒清洁工具设计；
\item 大容量强磁吸附工具设计；
\item 多功能过滤工具设计；
\item 可变径刮管工具设计；
\item “三合一”脉冲式冲刮刷集成式井筒清洁工具设计。
\end{circlist}

\subsection{工具应用效果的多因素分析方法研究}
\begin{circlist}
\item 收集工具的关键样本数据，构建工具“设计—应用”数据集；
\item 基于相关性分析，探究影响工具应用效果的主控因素。
\end{circlist}

\subsection{井筒清洁工具性能评价与作业参数优选研究}
\begin{circlist}
\item 基于实验系统，设计评价工艺，优选测试评价方法，开展“数据 + 机理”融合的井筒清洁工具性能评价研究，分析评估井筒清洁关键工具功能、耐用性、作业效率等使役性能，寿命等；
\item 基于“设计—应用”数据集，建立以作业效率为目标函数，以钻压、转速和排量、附着物、温度等为变量的大数据模型，考虑多参数之间的相互作用和影响，结合优化算法，形成作业参数智能优化方法，优选工具最佳工艺及作业参数组合。
\end{circlist}

\subsection{评价实验过程运转映射数字孪生技术研究}
\begin{circlist}
\item 工具评价实验过程的虚拟空间运转映射副本构建；
\item 监测工具评价实验过程，优化评价实验流程与方法。
\end{circlist}

\section{技术要求}

上述服务内容的技术要求包括但不限于以下内容：

\subsection{深度清洁工艺技术现状及适应性研究}
调研国内外先进的井筒清洁工艺及工具，对比分析各自的优缺点，结合海上油田现场作业条件，预测现场应用前景，形成详细的研究文件。

\subsection{井筒高效清洁技术配套理论研究}
\begin{enumerate}[label=(\arabic*),leftmargin=2em,itemsep=0pt]
\item 运用动力学、静力学分析，模拟砂砾、铁屑、结垢物、蜡等在井筒内不同流速、不同粘度等环境下的运移及分布规律研究，建立相关模型并基于室内试验验证规律的合理性；
\item 以井筒环境为应用基础，调研高强磁性体及高温磁性体的性能，对比筛选出最优产品，建立数值模拟模型，优化磁性体最大磁性布置规律、给出磁性体安放位置建议，结合磁性体消磁因素分析，给出最优的应用方式，提高工具使用寿命；
\item 建立水力脉冲模型，通过数值模拟或其他研究方式分析水力脉冲对井壁清理和携屑能力的影响，分析脉冲频率对井壁清理和携屑能力的分布规律，给出详细的分析研究文件；
\item 模拟刮管器的刮板在不同弹力的作用下对井壁附着物的清洁效果影响，以井筒环境为应用基础，在不影响管柱起下的前提下给出最优弹力分布范围；模拟不同钢刷材质和钢丝长度对井壁附着物的清洁效果影响，给出在不影响管柱起下的前提下最优材质和长度分布范围。
\end{enumerate}

\subsection{井筒高效清洁技术配套工具设计}
通过数值模拟和力学分析，对下列工具和关键组件进行结构优化、强度优化、组合优化和功能优化，绘制工具图纸并满足但不限于下列技术参数要求：
\begin{enumerate}[label=(\arabic*),leftmargin=2em,itemsep=0pt]
\item 四合一集成式井筒清洁工具，具备通刮刷磁吸功能，磁性材料耐温 $\ge 150^\circ\text{C}$，抗拉 $\ge 100\,\text{T}$，抗扭 $>17\ \text{kN}\cdot\text{m}$；
\item 大容量强磁吸附工具：抗拉 $\ge 100\,\text{T}$，抗扭 $\ge 18\ \text{kN}\cdot\text{m}$，吸附力 $\ge 8\ \text{kg/cm}^2$，吸附容量 $\ge 50\ \text{L}$；
\item 多功能过滤工具：工作允许旋转 $\ge 120\ \text{rpm}$，工具内通径 $\ge 2.688''$，碎屑存储量 $\ge 30\,\text{L}$；
\item 可变径刮管工具：可一趟钻完成 $9\frac{5}{8}''$ 套管和 $7''$ 套管刮管；
\item “三合一”集成式井筒清洁工具，具备脉冲、防卡、可变径功能，实现冲刮刷清洁井筒，可通过 $4.75''$ 密封筒、清洁 $9\frac{1}{2}''$ 筛盲管，耐温 $\ge 150^\circ\text{C}$，抗拉 $\ge 40\,\text{T}$，抗扭 $\ge 8\ \text{kN}\cdot\text{m}$。
\end{enumerate}

\subsection{工具应用效果的多因素分析方法研究}
\begin{enumerate}[label=(\arabic*),leftmargin=2em,itemsep=0pt]
\item 构建数据集 (工具$\ge 2$ 种)，数据集包含结构参数、材料参数、使役工况参数和性能指标参数等多维信息；
\item 开展工具应用效果主控因素识别研究 (工具$\ge 2$ 种)，通过相关性分析，确定对工具性能影响最为显著的关键参数。
\end{enumerate}

\subsection{井筒清洁工具性能评价与作业参数优选研究}
\begin{enumerate}[label=(\arabic*),leftmargin=2em,itemsep=0pt]
\item 评价 $9\frac{5}{8}''$ 井筒深度清洁研发的关键工具的功能、耐用性、作业效率等使役性能，形成井筒清洁工具应用效果室内测试评价工艺；
\item 运用优化算法对工具进行整体性能参数优选，形成作业参数智能优化方法。
\end{enumerate}

\subsection{评价实验过程运转映射数字孪生技术研究}
\begin{enumerate}[label=(\arabic*),leftmargin=2em,itemsep=0pt]
\item 采用数字孪生技术，在虚拟空间中实现工具评价实验过程的运转映射；
\item 工具工作轨迹模拟与实际作业轨迹的重合度达到 $90\%$ 以上。
\end{enumerate}

\subsection{乙方向甲方提交成果}
\begin{enumerate}[label=(\arabic*),leftmargin=2em,itemsep=0pt]
\item 提交项目成果总结报告、分析测试 (实验) 报告及原始数据；
\item 形成一套 $9\frac{5}{8}''$ 井筒深度清洁工艺技术；
\item 建立一套修井关键工具“设计—应用”数据集及实验评价方法；
\item 提供井筒清洁关键工具设计图纸各一套 (“四合一”集成式工具、大容量磁铁、多功能过滤器工具规格尺寸型号$\ge 3$ 种)；
\item 撰写专利 3 件，其中发明专利至少 1 件；
\item 录用或发表核心及以上期刊或 SCI/EI 论文至少 2 篇。
\end{enumerate}

\subsection{知识产权保护}
\begin{enumerate}[label=(\arabic*),leftmargin=2em,itemsep=0pt]
\item 项目中的成果知识产权归甲方所有；
\item 成果不侵犯任何第三方的合法权益，包括但不限于任何第三方的商业秘密、商标、服务标记、著作权、专利或其他任何知识产权。
\end{enumerate}